% sections/14_conclusion.tex

\section{Conclusion}

This report used a hypothesis-driven, visualisation-centric approach to analyse
the US Accidents dataset (2016--2022). The analysis spans geographic mapping,
temporal trend analysis, exploratory distribution analysis (box plots),
weather comparison, road feature investigation, and hourly severity
decomposition --- all grounded in a cleaned, $2\times$IQR outlier-removed
dataset that excludes the partial 2023 year.

\begin{table}[H]
\centering\small
\begin{tabularx}{\linewidth}{@{}lll X@{}}
  \toprule
  \hdrrow \textbf{Hypothesis} & \textbf{Verdict} & \textbf{Key Numbers} & \textbf{Key Insight} \\ \midrule
  H1: Bad weather
    & Partially Supported
    & Rain 24.38\%, Clear 17.85\%
    & Rain--Clear gap widened to 6.53\,pp after cleaning; Fog/Haze still underperforms. \\[4pt]
  H2: Rush hours
    & Partially Supported
    & 7--8\,am rate: 18.8--19\%; 3--6\,am plateau: 21.1--21.7\%
    & Rush hours have highest absolute counts, not highest rates. \\[4pt]
  H3: Infrastructure vs weather
    & Supported
    & Signal: $-11.7$\,pp; Rain vs Clear: $+6.53$\,pp
    & Infrastructure effect is $\approx1.8\times$ the weather effect. \\[4pt]
  H4: Low-traffic severity
    & Supported
    & Sat 7--9\,am peak $\approx$0.3176; visibility $>$10\,mi = 31.65\%
    & Sustained 3--6\,am plateau confirms speed/response-time mechanism. \\
  \bottomrule
\end{tabularx}
\caption{Summary of hypothesis verdicts and key insights (v3, post-cleaning)}
\end{table}

A critical finding across all hypotheses is that severity in this dataset
measures \emph{traffic disruption duration}, not physical harm. This reframes
the problem:

\begin{itemize}[noitemsep]
  \item \textbf{Infrastructure investment} (signals, crossings) is the most
        powerful lever for reducing high-severity incidents --- more so than
        weather-driven interventions.
  \item \textbf{Rush hours} matter for absolute count management (more
        resources, faster dispatch needed), but not for per-incident severity
        reduction.
  \item \textbf{Late-night and early morning hours} demand attention for
        per-incident severity reduction --- speed reduction measures and improved
        rural emergency response times are the targeted interventions.
  \item \textbf{Outlier removal materially improved the analysis:} adverse
        weather severity rates are higher, the early morning plateau is cleaner,
        and the Snow/Ice signal is stronger in the cleaned dataset.
\end{itemize}
